\documentclass[11pt,a4paper]{article}
\usepackage[utf8]{inputenc}
\usepackage[T1]{fontenc}
\usepackage[french]{babel}
\usepackage{amsmath,amssymb,amsthm,geometry,hyperref}
\geometry{margin=2.4cm}
\newtheorem{proposition}{Proposition}
\title{Changement de support auxiliaire avec mémoire FE\\
Transition contractante et mesure de la mémoire supprimée}
\author{Étude GPU SLFE, fork keckj/rodin}
\date{16 novembre 2026}
\begin{document}
\maketitle

\section{Le problème du support variable}
Le couplage avec mémoire possède une preuve d'énergie pour un transfert
fixé $R:V\to W$. Changer le support $W$ et remplacer simplement $R$ dans
les mises à jour ne conserve pas cette preuve. Le diagnostic étudié
ici maintient le champ FE inchangé au changement et reconstruit un état
augmenté compatible avec le nouveau support.

L'espace FE $V$ est fixé. Chaque support $W_k$ est muni de sa norme de
masse et d'une contraction $R_k:V\to W_k$ préservant les constantes,
de même que son adjoint. On pose
\[
 D_k=(I-R_k^*R_k)^{1/2},\qquad J_k v=(R_kv,D_kv).
\]
Les $J_k$ sont des isométries. Entre changements, le calcul conserve
trois tableaux $v,w,a$ vérifiant
\[
 Z=(w,D_ka),\qquad v=J_k^*Z=R_k^*w+(I-R_k^*R_k)a.
\]
L'initialisation est $a=v_0$, $w=R_kv_0$. Une advection modifie uniquement
$w$ dans l'espace augmenté. Une opération physique FE de variation
$\delta v$ ajoute $R_k\delta v$ à $w$ et $\delta v$ à $a$.
Ces mises à jour conservent l'identité précédente.

\section{Une énergie accessible sans racine de matrice}
L'énergie augmentée est calculable directement par
\[
 E=\|w\|_{W_k}^2+\|a\|_V^2-\|R_ka\|_{W_k}^2.
\]
Cette formule ne requiert ni $D_k$, ni résolution linéaire. Les opérations
d'advection et de physique considérées restent contractantes sous les
hypothèses établies pour le couplage à support fixé.

\begin{proposition}
La remise à zéro $a^+=v$, $w^+=R_{k+1}v$ préserve exactement les coefficients
FE au changement de support. La transition augmentée est
$Z^+=J_{k+1}J_k^*Z$ et elle est contractante. Sa perte d'énergie est
\[
 E-E^+=E-\|v\|_V^2=\|(I-J_kJ_k^*)Z\|^2\geq0.
\]
\end{proposition}
\begin{proof}
Le champ observé après transition est
$J_{k+1}^*Z^+=J_{k+1}^*J_{k+1}J_k^*Z=v$.
L'isométrie de $J_{k+1}$ donne $\|Z^+\|^2=\|v\|^2$.
L'opérateur $J_kJ_k^*$ est un projecteur orthogonal; l'identité de
Pythagore donne l'expression de la perte.
\end{proof}

La quantité supprimée mesure une composante de mémoire invisible dans
l'observation FE au moment du changement. Ce n'est pas une borne d'erreur
par rapport à l'EDP: le transport futur peut rendre cette composante
observable, et le nouveau support introduit ses propres erreurs spatiales.
Des pertes légèrement négatives peuvent apparaître par soustraction en
virgule flottante; les journaux les conservent sans troncature.

Les constantes et la masse FE sont préservées au changement. Toute suite
de transports et de physiques contractants, séparée par ces transitions,
reste bornée par l'énergie initiale. Cette propriété reste vraie si la
décision de changer de support dépend du champ calculé: chaque transition
choisie conserve individuellement la même inégalité. Aucune linéarité du
contrôleur d'adaptation n'est requise.

\section{Transférer une partie de la mémoire invisible}
Soit $C:W_0\to W_1$ une contraction entre les normes auxiliaires.
Avec $Q_k=I-J_kJ_k^*$, la transition
\[
 B=J_1J_0^*+Q_1\begin{pmatrix}C&0\\0&0\end{pmatrix}Q_0
\]
préserve l'observation FE et est contractante. Les deux termes ont des
images orthogonales et agissent sur les composantes orthogonales
$J_0J_0^*Z$ et $Q_0Z$, respectivement.

Cette formule possède une réalisation sans racine de matrice. Posons
\[
 r=w-R_0v,\qquad h=Cr,\qquad b=R_1^*h,
 \qquad a^+=v-b,\qquad w^+=R_1a^++h.
\]
Le champ $v$ reste inchangé, puisque
$R_1^*w^++(I-R_1^*R_1)a^+=v$. La nouvelle énergie cachée vaut
\[
 E^+-\|v\|^2=\|h\|^2-\|R_1^*h\|^2.
\]
La perte possède la décomposition exacte
\[
 E-E^+=\underbrace{\|D_0(a-v)\|^2}_{\text{mémoire du défaut FE}}
       +\underbrace{\|r\|^2-\|Cr\|^2}_{\text{projection entre supports}}
       +\underbrace{\|R_1^*Cr\|^2}_{\text{correction de l'observation}}.
\]
Les trois termes sont non négatifs. Le premier est calculable par
$\|a-v\|^2-\|R_0(a-v)\|^2$. La remise à zéro complète correspond à
$C=0$. Pour tout $C$ contractant, la perte est inférieure ou égale à
celle de cette remise à zéro, pour un même état initial. Cette comparaison
locale des énergies n'implique pas une réduction de l'erreur globale de
transport; cette dernière est mesurée séparément.

Les validations denses GPU aux degrés 1, 2 et 4 reconstruisent $D_k$,
$J_k$, $Q_k$ et $B$ indépendamment des mises à jour à trois tableaux.
Elles contrôlent la norme de $B$, l'identité $J_1^*B=J_0^*$ et les trois
termes de perte. Sur le transport cellulaire testé avec six changements,
le transfert de mémoire réduit l'erreur $L^2$ de $0{,}00620$ à $0{,}00385$
sur le maillage FE $64^2P_4$, dont le plancher est $0{,}00373$.
Sur $128^2P_4$, les erreurs correspondantes sont $1{,}98\,10^{-4}$ et
$1{,}60\,10^{-4}$, pour un plancher de $1{,}43\,10^{-4}$.

\section{Raffinement emboîté sans perte de mémoire}
Un cas structurel mérite d'être distingué. Supposons que les supports
soient des sous-espaces du même $L^2$ et que $W_0\subset W_1$.
L'injection isométrique $C:W_0\to W_1$ vérifie $C^*R_1=R_0$.
Alors la mise à jour
\[
 a^+=a,\qquad w^+=Cw+R_1a-CR_0a
\]
conserve exactement $v$ et l'énergie augmentée. En effet, dans la
représentation ambiante $q=w+(I-P_{W_0})a$, la nouvelle décomposition
est $q=w^++(I-P_{W_1})a$; le champ ambiant n'a pas changé.
La conservation de la norme suit de l'orthogonalité des deux composantes.
La formule injecte aussi les modes de l'ancre FE nouvellement résolus.
Ce résultat suppose un emboîtement exact. Il ne s'applique pas directement
au déplacement non emboîté des partitions graduées.
L'implémentation GPU est contrôlée aux degrés 1, 2 et 4, pour un
raffinement global et un raffinement local. Les matrices de Gram de
l'état augmenté avant et après changement diffèrent de moins de
$4{,}3\,10^{-15}$ sur ces six cas; l'identité d'observation FE diffère
de moins de $3{,}4\,10^{-15}$. Ces contrôles d'algèbre linéaire ne valident
pas à eux seuls un critère de choix des cellules à raffiner.

\section{Projection itérative dans un espace ambiant commun}
Une transition plus informative est possible lorsque $R_k=P_{W_k}|_V$
et $C=P_{W_1}|_{W_0}$ sont les projections exactes du même espace $L^2$.
Une contraction $C$ abstraite ne suffit plus. L'état possède alors la
représentation physique
\[
 q=w+(I-P_{W_0})a,\qquad v=P_Vq,\qquad \|q\|^2=E.
\]
Posons $H=(I-P_V)|_{W_1}$ et $K=H^*H=I-R_1R_1^*$.
La donnée pour projeter la mémoire invisible est calculable sans
reconstruire $q$ :
\[
 h=H^*(q-v)=C(w-R_0a)+R_1(a-v).
\]
Partant de $x_0=0$, on effectue $k$ itérations
\[
 x_{j+1}=h+R_1R_1^*x_j,
 \qquad a^+=v-R_1^*x_k,\qquad w^+=R_1a^++x_k.
\]
Le nouvel état ambiant est $q^+=v+Hx_k$, donc son observation FE
reste exactement $v$.

\begin{proposition}
Pour chaque entier $k\geq1$, cette transition est contractante.
Pour un état entrant fixé, l'énergie de mémoire conservée croît avec
$k$, et $q^+$ converge vers $P_{V+W_1}q$ lorsque $k\to\infty$.
\end{proposition}
\begin{proof}
Les valeurs propres de $K$ appartiennent à $[0,1]$. Sur un vecteur
singulier gauche de $H$ associé à $\lambda>0$, l'opérateur
\[
 H\sum_{j=0}^{k-1}(I-K)^jH^*
\]
multiplie la mémoire par $g_k(\lambda)=1-(1-\lambda)^k$.
Ce facteur appartient à $[0,1]$, croît avec $k$, et tend vers 1.
Sur l'orthogonal de l'image de $H$, le facteur est nul.
La partie dans $V$ est conservée et orthogonale à ces composantes.
Enfin $V+W_1=V\mathbin{\oplus}\operatorname{im}H$, ce qui identifie
la limite au projecteur orthogonal annoncé.
\end{proof}

Cette transition distingue deux notions : conserver la mémoire
approximable sur le nouveau support et diminuer l'erreur par rapport à
l'EDP. La première propriété est démontrée; la seconde reste une question
expérimentale. La convergence de l'itération peut être lente lorsque
$K$ possède de petites valeurs propres non nulles, c'est-à-dire lorsque
les espaces présentent de petits angles principaux. Les itérations
finies ne requièrent ni racine, ni inverse, ni estimation du noyau de $K$.

L'oracle GPU reconstruit indépendamment le Gram conjoint de $V,W_0,W_1$,
plonge ces espaces dans un même espace euclidien et applique le filtre
par diagonalisation. Les degrés 1, 2 et 4 et les nombres d'itérations
1, 2, 4, 8, 16 et 32 sont contrôlés. Les écarts d'état sont inférieurs
à $10^{-14}$ et les écarts d'énergie absolus à $2{,}3\,10^{-13}$ pour
des énergies initiales de l'ordre de plusieurs centaines.

\section{Choix du support et limites de coût}
Le premier contrôleur utilise le centre circulaire de l'énergie du champ,
calculé aux points de Gauss. Il déplace le minimum de taille des cellules
lorsque ce centre s'écarte d'une distance prescrite. Il s'agit d'une
heuristique pour une structure localisée unique; le module des moments
circulaires est enregistré pour détecter une perte de concentration.
Les cas à plusieurs structures séparées demandent un autre choix de
support. La reconstruction FE--auxiliaire au changement est incluse
dans les coûts mesurés.

Le transport du seul état DG entre supports mobiles est également une
référence nécessaire. Les projections tensorielles successives dans les
deux directions préservent masse et contraction. Cette variante peut
éviter les reconstructions FE intermédiaires pour l'advection pure, au
prix de l'abandon du défaut FE initial. La comparaison distingue ce
choix du couplage avec mémoire et physique FE.

\section{Antécédents et revendications}
L'itération de projection et son interprétation par les angles entre
sous-espaces appartiennent à une théorie classique; l'itération de
Richardson ci-dessus n'est pas revendiquée comme un nouvel algorithme
d'algèbre linéaire. Bauschke et Borwein exposent le cadre des projections
alternées \href{https://cmps-people.ok.ubc.ca/bauschke/Research/01.pdf}{(article)}.
Adams, Finn, Knepley et Pusztay étudient déjà le rééchantillonnage de
particules par projection FE et pseudo-inverse, avec conservation des
moments représentables
\href{https://arxiv.org/abs/2501.13681}{(prépublication 2025, publication 2026)}.
La question à comparer précisément est ici le maintien du champ FE et
de sa mémoire augmentée à chaque troncature de l'itération.

L'adaptation SL-DG est établie, notamment par Besse, Deriaz et Madaule,
\emph{Adaptive multiresolution semi-Lagrangian discontinuous Galerkin
methods for the Vlasov equations}, JCP 332 (2017), 376--417
\href{https://www.oca.eu/images/LAGRANGE/pages_perso/nbesse/pub/NBesse_jcp17.pdf}{(article)}.
Adams étend les intégrales de recouvrement DG à des grilles AMR
\href{https://arxiv.org/abs/2603.19959}{(prépublication, 2026)}.
Le changement de taille des intervalles et l'adaptation ne sont donc
pas des revendications de nouveauté.

CO-FLIP utilise des transferts par projection et une remise à zéro
adaptative des cartes de flot, déclenchée par leur déformation
\href{https://cseweb.ucsd.edu/~ravir/sinasiga.pdf}{(Nabizadeh et al., 2024)}.
La remise à zéro étudiée ici concerne un état augmenté FE--auxiliaire
linéaire et son défaut de projection. La proximité conceptuelle avec
ces antécédents impose une analyse de nouveauté au-delà de l'identité
élémentaire de projection démontrée ci-dessus.
\end{document}
